\section{Themes and Determinism}
\label{sec:themes-determinism}

\subsection{Themes}

Triton separates meaning from appearance. The grammar fixes what a diagram
\emph{is}; the theme fixes how it \emph{looks}. A diagram never hard-codes a
colour, and a theme never changes a diagram's structure.

A theme is one resolved token set (\texttt{src/contracts/theme.ts}) ---
\texttt{ResolvedTheme} --- with token groups for palette, typography, spacing,
edges, and panel chrome:

\begin{lstlisting}[language=ts]
interface ResolvedTheme {
  palette:    ThemePalette;     // primary, surface, border, text, success, ...
  typography: ThemeTypography;  // font families, sizes, line height
  spacing:    ThemeSpacing;
  edges:      ThemeEdges;
  panel:      ThemePanel;       // titled-container chrome
}
\end{lstlisting}

A resolver merges a user override over a base preset; the codebase ships a dozen
presets (\texttt{src/theme/preset.ts}). One resolved theme applies uniformly to
\emph{every} diagram kind --- switching themes restyles a flowchart, a tree, and a
poster identically, because they all read the same tokens through the
\texttt{Pen}.

\subsection{Determinism}

The same source and theme always produce the same SVG, byte for byte. This is a
property the whole pipeline is built to preserve, not an afterthought:

\begin{itemize}
  \item Layout is a pure function of the IR and theme --- no clocks, no random
        numbers, no hash-order iteration over unordered maps.
  \item Text measurement (\texttt{measureText}) is deterministic, so geometry does
        not depend on a live font engine.
  \item Numeric output is formatted at fixed precision, and elements are emitted
        in a stable document order.
  \item The value-driven builders are deterministic by construction: the same
        \texttt{insert} list yields the same tree.
\end{itemize}

Byte-stable SVG is what makes the output diff-able, cacheable, and safe to commit
beside its source --- which is exactly how the \texttt{examples/} directory is
organised, each \texttt{.svg} regenerated from its \texttt{.mmd}.
